\section{Electromagnetic potentials in quantum mechanics}

Up to this point nothing has forced us to take the potentials seriously. The
classical equations of motion contain only \(E\) and \(B\), and the gauge
freedom
\[
	A \longmapsto A + \nabla \chi, \qquad V \longmapsto V -
	\frac{\partial \chi}{\partial t},
\]
with \(\chi(r,t)\) an arbitrary function, never reaches an observable.
Quantum mechanics removes this comfort.

\subsection{Mathematical Interlude II: The Quantum Hamiltonian}

In classical mechanics, the motion of a particle is determined by the force
acting on it: Newton's seconds law fixes the acceleration, and the
trajectory follows by integration.
Quantum mechanics is formulated differently. The state of the particle
is described by a wavefunction \(\psi(r,t)\), whose time evolution is
governed by an operator called the Hamiltonian,
\[
	i \hbar \frac{\partial \psi}{\partial t} = \hat H \psi
	\label{eq:schrodinger}
\]
For the non-relativistic system considered here, the Hamiltonian
represents the total energy of the particle.



The Hamiltonian is the operator associated with the total energy of the
system. The \cref{eq:schrodinger}
states that the energy operator \emph{generates} the evolution:
the change of \(\psi\) during an infinitesimal interval \(dt\) is
prescribed by \(\hat H \psi\). To specify the dynamics of a quantum system is
therefore to specify its energy, in the same sense in which specifying the
dynamics of a classical particle means specifying the force acting on it.


Classical mechanics admits an equivalent formulation, due to Hamilton, in
which the evolution is generated by the function
\vspace{-5pt}
\[
	H (r, p) = \frac{p^2}{2m}+ U(r), \qquad \dot r =
	\frac{\partial H}{\partial p}, \qquad
	\dot p = - \frac{\partial H}{\partial r},
\]
rather than by \(F\). The trajectories obtained are exactly the Newtonian ones;
what changes its the object taken as primitive. The Hamilton's formulation is
written in terms of the potential energy, rather than Newton's that is
written in forces. In quantum mechanics this choice becomes a choice of
what the theory regards as fundamental.


The simplest consequence of \cref{eq:schrodinger}
fixes the meaning of the Hamiltonian. If \(U\) does not depend on time,
there exist solutions of definite energy,
\vspace{-3pt}
\[
	\psi(r,t) = \psi_E(r)\, e^{-iEt \hbar},
\]
for which the probability density \(|\psi|^2 = |\psi_E|^2\) is independent of
time. The time evolution of a stationary state consist entirely of a phase
factor. This is the first indication of a structural feature of quantum
mechanics: a global phase, \(\psi \longmapsto e^{i \alpha}\psi\) with
\(\alpha\) constant, leaves every prediction unchanged and is therefore not
an observable. A \emph{relative} phase between two components of a
superposition is a different matter.

\enlargethispage{0.7cm}
If \(\psi = \psi_1 + \psi_2\) then
\[
	|\psi|^2 = |\psi_1|^2 + |\psi_2|^2 + 2 \Re \!
	\left(
	\psi_1^* \psi_2
	\right)
\]
and the interference term depends on the phase difference between the two
branches. Quantum mechanics thus provides an observable that has no classical
counterpart, and it is though this observable that the electromagnetic
potentials will enter the physics.


\section{Minimal coupling}

A quick recap of the basics of quantum mechanics is in 
order before proceeding.

In quantum mechanics the state of a particle is a function 
\(\psi \colon \mathbb{R} \to \mathbb{C}\). In every point of the space 
we have a complex number; this can be visualized as a bizarre clock. The norm
\(|\psi|\) tell us the probability of finding the particle in that precise 
spot in space; the argument tell us where the lancet is pointing.

\begin{figure}[H]
	\centering
	\begin{tikzpicture}[
	line cap=round, line join=round,
	axisline/.style={-{Latex[length=2mm]}, line width=0.4pt, black!80},
	dial/.style={line width=0.35pt, black!45, dash pattern=on 1.1pt off 1.3pt},
	hand/.style={-{Latex[length=1.4mm,width=1.1mm]}, line width=0.7pt, black},
	guide/.style={line width=0.3pt, black!35, dotted},
	curve/.style={line width=0.8pt, black},
	lbl/.style={font=\footnotesize}
	]

	%================ (a) the value of psi at one point ================
	\begin{scope}
		\def\R{1.40}
		\pgfmathsetmacro{\th}{54}
		\pgfmathsetmacro{\rr}{1.05}

		\draw[axisline] (-\R-0.30,0) -- (\R+0.50,0) node[right,lbl] {$\operatorname{Re}$};
		\draw[axisline] (0,-\R-0.30) -- (0,\R+0.50) node[above,lbl] {$\operatorname{Im}$};

		\draw[dial] (0,0) circle (\R);

		\draw[guide] (\th:\rr) -- ({\rr*cos(\th)},0);
		\draw[guide] (\th:\rr) -- (0,{\rr*sin(\th)});

		\draw[line width=0.4pt, black!80] (0.40,0) arc (0:\th:0.40);
		\node[lbl] at ({0.68*cos(0.5*\th)},{0.68*sin(0.5*\th)}) {$\theta$};

		\draw[hand] (0,0) -- (\th:\rr);
		\fill (0,0) circle (0.035);

		\path (0,0) -- (\th:\rr)
		node[midway, sloped, above, lbl, inner sep=1.5pt] {$|\psi|$};

		\node[lbl] at (0,-\R-0.75) {$\psi(x)=|\psi|\,e^{i\theta}$};
	\end{scope}

	%================ (b) psi along space ================
	\begin{scope}[shift={(-4.05,-4.65)}]
		\def\dx{1.35}
		\def\r{0.52}
		\def\yc{-1.55}
		\def\hmax{1.20}
		\pgfmathsetmacro{\xc}{3*\dx}
		\pgfmathsetmacro{\sig}{2.6*\dx}

		% modulus profile
		\draw[curve] plot[domain=-0.85:8.95, samples=180, smooth]
		(\x,{\hmax*exp(-0.5*(\x-\xc)*(\x-\xc)/(\sig*\sig))});
		\node[lbl, above right] at (7.05,{\hmax*exp(-0.5*(7.05-\xc)*(7.05-\xc)/(\sig*\sig))})
		{$|\psi(x)|$};

		% axis
		\draw[axisline] (-1.05,0) -- (9.15,0) node[right,lbl] {$x$};

		\foreach \i in {0,...,6}{
				\pgfmathsetmacro{\xi}{\i*\dx}
				\pgfmathsetmacro{\amp}{exp(-0.5*(\xi-\xc)*(\xi-\xc)/(\sig*\sig))}
				\pgfmathsetmacro{\ang}{\i*45}
				\draw[line width=0.4pt, black!80] (\xi,0) -- (\xi,-0.10);
				\draw[guide] (\xi,{\hmax*\amp}) -- (\xi,{\yc+\r+0.12});
				\draw[dial] (\xi,\yc) circle (\r);
				\draw[hand] (\xi,\yc) -- ++(\ang:{\r*\amp});
				\fill (\xi,\yc) circle (0.03);
			}

		\node[lbl] at (\xc,{\yc-\r-0.55}) {$\psi:\mathbb{R}^3\longrightarrow\mathbb{C}$};
	\end{scope}

\end{tikzpicture}

	\caption{Clock representation of a complex wavefunction. At each position,
		the length of the hand gives the modulus \(\abs{\psi(x)}\), while its angle
		gives the phase of \(\psi(x)\).}
	\label{fig:wavefunction-clock}
\end{figure}

The norm is measurable, the lancet is not. If we rotate

